\documentclass[11pt,a4paper]{article}
\usepackage[utf8]{inputenc}
\usepackage[margin=0.9in]{geometry}
\usepackage{parskip}
\usepackage{setspace}
\pagenumbering{gobble}
\setstretch{1.1}
% Harvard Version
\begin{document}

% \begin{center}
% \textbf{Personal Statement}
% \end{center}

\begin{center}
\Large\textbf{Personal Statement}\\[4pt]
\normalsize Benedikt Samwer\quad|\quad PhD Applicant, Biomedical Informatics \quad|\quad Harvard University
\end{center}

After finishing school I spent a year as a full time volunteer in a major hospital in Haifa (Israel), working on a pediatric ward. I helped with small but essential tasks for patients and families, translated between staff and international visitors, and watched nurses and doctors juggle enormous responsibility with limited time and resources. It was my first close view of clinical practice. I loved being around patients and seeing care teams at work, but I also noticed how much depended on individual memory, paper forms, and fragmented systems. That year convinced me that I wanted my work to serve patients and clinicians, but also made me aware that my strongest skills were quantitative rather than clinical.

Back in Germany I chose Engineering Science at the Technical University of Munich because it let me go deep into mathematics, data analysis, and modeling while keeping an eye on real systems. For my bachelor thesis at the Research and Consulting Center for Energy Economics I worked on simulation and optimization of district heating networks and storage systems. I learned to frame large constrained optimization problems, but also to sit in workshops with city planners and explain model behavior in plain language. The experience of moving between equations and real world constraints felt very similar to conversations I had watched in the pediatric ward between specialists and families. It reinforced that I enjoy being in the middle, translating between domains.

The desire to work closer to medicine never really went away, so I enrolled in Medicine at LMU for one year in parallel to my engineering studies. Studying anatomy and biochemistry, and shadowing on wards, made it clear how much pressure clinicians are under and how many decisions they make with incomplete information. It sharpened my sense that I could contribute more by building systems that support them than by becoming a clinician myself. Seeing both the human side of care in Haifa and Munich and the structural limits of current workflows is what pushed me toward AI in medicine as a way to combine my academic strengths with my wish to work for the public good.

Community building has become a natural extension of this bridge role. In 2024 I co founded “med-dev,” a student and young professional community that brings clinicians and machine learning practitioners together. We organize talks and small hack sessions on reproducible ML and clear communication, with a focus on making the space welcoming for people who feel outside their comfort zone, whether that is clinicians trying to understand models or tech students seeing a hospital environment for the first time. Designing these events has taught me about facilitation, psychological safety, and the value of informal peer networks.

Teaching has been another constant. As a tutor and teaching assistant at TUM I led mechanics and mathematics tutorials for more than 100 students, many of whom came in convinced they were “not good at math.” I learned to break problems into small, honest steps and to give feedback that built confidence instead of just displaying solutions. That patience and focus on clarity now shape how I mentor peers in research projects.

Having lived and worked in Germany, Israel, the United Kingdom, and Boston, I have become comfortable moving between cultures and disciplines. At Harvard Medical, within the Biomedical Informatics PhDs AIM track, I hope to contribute not only through my research, but also as someone who builds bridges between engineers and clinicians, international and local students, and theory and practice.


\end{document}
